\indent ಬ್ರಹ್ಮಸೂತ್ರ : ಸಂಧಿಕಾರ್ಯ ನಡೆಯುವಾಗ ಅ, ಆ ಕಾರಗಳಲ್ಲಿ ಯಾವುದಾದರೊಂದು ಸ್ವರಕ್ಕೆ ಏ, ಐ ಕಾರಗಳು ಪರವಾದರೆ ಐ ಕಾರವೂ ಓ, ಔ ಕಾರಗಳು ಪರವಾದರೆ ಔಕಾರವೂ ಆದೇಶಗಳಾಗಿ ಬರುವುದಾದರೆ ಅದು ವೃದ್ಧಿಸಂಧಿಯೆನಿಸುವುದು.\par

\vspace{4pt plus 4pt minus 1pt}\noindent{\kbvsecfont\fontsize{11pt}{14.5pt}\selectfont\color{kbvaccentdark} {[}೪{]}\hspace{0.45em}ಯಣ್‌ ಸಂಧಿ :}\hspace{0.5em}ಇತ್ಯಾದಿ, ಗುರ್ವಾಜ್ಞೆ, ನಿತ್ರಾರ್ಜಿತ - ಈ ಶಬ್ದಗಳನ್ನು ಬಿಡಿಸಿ ನೋಡೋಣ...\par

\begin{kbvderiv}
\vspace{4pt plus 3pt minus 1pt}
\begin{center}\begin{adjustbox}{max width=\textwidth}
{\renewcommand{\arraystretch}{1.22}\setlength{\tabcolsep}{4pt}
\begin{tabular}{l c l}
{[}{\kbvlatin i}{]} ಇತ್ಯಾದಿ & {\color{kbvgrey}$=$} & ಇತಿ + ಆದಿ \\
 & {\color{kbvgrey}$=$} & ಇ + {[}ತ್ + ಆ + ದಿ{]} + ಆ \\
\end{tabular}}
\end{adjustbox}\end{center}
\vspace{4pt plus 3pt minus 1pt}
\end{kbvderiv}

\indent ಸ ಇಕಾರಕ್ಕೆ ಯ್‌ಕಾರಾದೇಶವಾಗಲು\par

\begin{kbvderiv}
\vspace{4pt plus 3pt minus 1pt}
\begin{center}\begin{adjustbox}{max width=\textwidth}
{\renewcommand{\arraystretch}{1.22}\setlength{\tabcolsep}{4pt}
\begin{tabular}{l c l}
 & {\color{kbvgrey}$=$} & ಇ + (ತ್ + ಆ + ದಿ) + ಆ \\
 & {\color{kbvgrey}$=$} & ಇತ್ಯಾದಿ.... {[}\ensuremath{\because} ತ್ + ಆ + ದಿ = ತ್ಯಾದಿ{]} \\
\end{tabular}}
\end{adjustbox}\end{center}
\vspace{4pt plus 3pt minus 1pt}
\end{kbvderiv}

\begin{kbvderiv}
\vspace{4pt plus 3pt minus 1pt}
\begin{center}\begin{adjustbox}{max width=\textwidth}
{\renewcommand{\arraystretch}{1.22}\setlength{\tabcolsep}{4pt}
\begin{tabular}{l c l}
{[}{\kbvlatin ii}{]} ಗುರ್ವಾಜ್ಞೆ & {\color{kbvgrey}$=$} & ಗುರು + ಆಜ್ಞೆ \\
 & {\color{kbvgrey}$=$} & ಗು + {[}ವ್ + ಆ + ಜ್ಞೆ{]} + ಆ \\
 & {\color{kbvgrey}$=$} & ಗುರ್ವಾಜ್ಞೆ.... {[}\ensuremath{\because} ವ್ + ಆ + ಜ್ಞೆ = ವಾಜ್ಞೆ{]} \\
\end{tabular}}
\end{adjustbox}\end{center}
\vspace{4pt plus 3pt minus 1pt}
\end{kbvderiv}

\begin{kbvderiv}
\vspace{4pt plus 3pt minus 1pt}
\begin{center}\begin{adjustbox}{max width=\textwidth}
{\renewcommand{\arraystretch}{1.22}\setlength{\tabcolsep}{4pt}
\begin{tabular}{l c l}
{[}{\kbvlatin iii}{]} ನಿತ್ರಾರ್ಜಿತ & {\color{kbvgrey}$=$} & ಪಿತೃ + ಆರ್ಜಿತ \\
 & {\color{kbvgrey}$=$} & ಪಿ + {[}ತ್ + ಋ + ಆ{]} + ರ್ಜಿತ \\
 & {\color{kbvgrey}$=$} & ಪಿ + (ತ್ + ಋ + ಆ) + ರ್ಜಿತ \\
 & {\color{kbvgrey}$=$} & ನಿತ್ರಾರ್ಜಿತ.... (\ensuremath{\because} ತ್ + ಋ + ಆ = ತ್ರಾ) \\
\end{tabular}}
\end{adjustbox}\end{center}
\vspace{4pt plus 3pt minus 1pt}
\end{kbvderiv}

\indent ಮೇಲಿನ ಉದಾಹರಣೆಗಳಲ್ಲಿ ಪೂರ್ವಾರ್ಧದಲ್ಲಿರುವ ಸ್ವರಾಕ್ಷರಗಳಾದ ಇ, ಉ, ಋ ಕಾರಗಳಿಗೆ ಕ್ರಮವಾಗಿ ಯ್‌, ವ್‌, ರ್‌ಗಳು ಆದೇಶವಾಗಿ ಬಂದಿರುವುದನ್ನು ನೋಡಬಹುದು ; ಇಲ್ಲಿ ಸವರ್ಣ ಸ್ವರ ಪರವಾಗಿ ಬಂದಿಲ್ಲದಿರುವುದು ಗಮನಾರ್ಹ ಅಂಶ. ಈ ತೆರನಾದ ಸಂಧಿಯು 'ಯಣ್‌ ಸಂಧಿ' ಎಂದು ಬಿಂದುಸುವುದು.\par
